\documentclass[a4paper,12pt]{article}
\usepackage[T2A]{fontenc}
\usepackage[utf8]{inputenc}
\usepackage[russian]{babel}
\usepackage{amsmath,amssymb,mathtools}
\usepackage{cmbright}
\renewcommand{\familydefault}{\sfdefault}
\usepackage{icomma}
\usepackage{geometry}
\geometry{left=23mm,right=23mm,top=23mm,bottom=23mm}
\usepackage{microtype}
\usepackage{tikz}
\usetikzlibrary{calc,arrows.meta,positioning,shapes.geometric}
\usepackage{booktabs,tabularx}
\usepackage{enumitem}
\usepackage{fancyhdr}
\usepackage{parskip}
\usepackage{setspace}
\usepackage{xcolor}

\definecolor{accent}{HTML}{176B70}
\definecolor{darkaccent}{HTML}{194F55}
\definecolor{textgray}{HTML}{303638}
\color{textgray}
\onehalfspacing
\setlength{\headheight}{16pt}
\pagestyle{fancy}
\fancyhf{}
\fancyhead[L]{\small Ксения Васильченко}
\fancyhead[R]{\small 15.10.26}
\fancyfoot[C]{\small\thepage}
\renewcommand{\headrulewidth}{0.35pt}
\setlist[enumerate]{leftmargin=*,itemsep=.9em,topsep=.35em}
\newcommand{\sectionline}[1]{\section*{\color{darkaccent}#1}}
\newcommand{\learn}[1]{\par\noindent$\square$\enspace #1\par}
\newcommand{\keyidea}[1]{\par\noindent\textcolor{darkaccent}{\textbf{Запомните.}}\enspace #1\par}
\newcommand{\mix}[3]{#1\,\frac{#2}{#3}}

\begin{document}
\begin{titlepage}
\thispagestyle{empty}
\centering
\vspace*{30mm}
{\color{accent}\large МАТЕМАТИКА \textbullet\ 6 КЛАСС\par}
\vspace{18mm}
{\color{darkaccent}\Huge\bfseries Чек-лист\par}
\vspace{9mm}
{\LARGE Дроби, уравнения\par}
\vspace{2mm}
{\LARGE и текстовые задачи\par}
\vspace{11mm}
{\large Удобные вычисления \textbullet\ НОК и НОД\par}
\vspace{2mm}
{\large Математический язык \textbullet\ Практические задачи\par}
\vspace{18mm}
{\large 15 октября 2026 года\par}
\vfill
{\large Лёвин Артём Александрович\par}
\vspace{3mm}
{\normalsize эксклюзивно для Ксении Васильченко\par}
\vspace{32mm}
\begin{tikzpicture}[remember picture,overlay]
\draw[accent!58,line width=1.25pt]
  ([xshift=25mm,yshift=20mm]current page.south west)
  .. controls ([xshift=79mm,yshift=43mm]current page.south west)
           and ([xshift=-77mm,yshift=5mm]current page.south east)
  .. ([xshift=-25mm,yshift=20mm]current page.south east);
\end{tikzpicture}
\end{titlepage}

\sectionline{1. Удобные вычисления с дробями}
\learn{Вы замечаете слагаемые с одинаковыми знаменателями и группируете их.}
\learn{Вы сохраняете все слагаемые на каждом промежуточном шаге.}

\textbf{Пример 1.} Сначала объединим дроби со знаменателем $15$:
\begin{align*}
\mix{2}{4}{15}+\mix{1}{7}{20}+\mix{3}{11}{15}
&=\left(\mix{2}{4}{15}+\mix{3}{11}{15}\right)+\mix{1}{7}{20}\\
&=6+\mix{1}{7}{20}
=\boxed{\mix{7}{7}{20}}.
\end{align*}
При \emph{сложении} слагаемые можно переставлять и группировать. При вычитании
важно сохранить знак, относящийся к каждому слагаемому.

\textbf{Пример 2.} Минус перед скобками меняет знаки внутри:
\[
a-(b+c)=a-b-c,\qquad a-(b-c)=a-b+c.
\]
Например,
\[
\mix{5}{1}{6}-\left(\mix{2}{1}{6}+\frac12\right)
=\mix{5}{1}{6}-\mix{2}{1}{6}-\frac12
=3-\frac12=\boxed{\mix{2}{1}{2}}.
\]
\keyidea{В длинном выражении записывайте и те части, до которых ещё не дошли.
Пропуск промежуточного слагаемого делает решение непоследовательным.}

\sectionline{2. Десятичные дроби и «заём» единицы}
\learn{Вы записываете десятичную дробь в виде смешанного числа и сокращаете её.}
\learn{Вы умеете занять одну единицу, если дробной части не хватает для вычитания.}

\[
1{,}9=\mix{1}{9}{10},
\qquad
2{,}34=\mix{2}{34}{100}=\mix{2}{17}{50}.
\]
Если дробные части имеют общий знаменатель $d$, то
\[
a+\frac{m}{d}=(a-1)+\frac{m+d}{d}
\qquad (a\geq1).
\]
Например, $10+\frac13=9+\frac43=8+\frac73$.
Значение числа остаётся прежним.

\textbf{Пример 3.} Переведём десятичную дробь в смешанное число:
\begin{align*}
\mix{2}{3}{15}-1{,}9+\mix{1}{7}{20}
&=\mix{2}{1}{5}-\mix{1}{9}{10}+\mix{1}{7}{20}\\
&=\mix{2}{4}{20}-\mix{1}{18}{20}+\mix{1}{7}{20}\\
&=\left(1+\frac{24}{20}\right)-\mix{1}{18}{20}+\mix{1}{7}{20}\\
&=\frac6{20}+\mix{1}{7}{20}
=\boxed{\mix{1}{13}{20}}.
\end{align*}
\keyidea{При «заёме» целая часть уменьшается на $1$, а числитель
дробной части увеличивается на знаменатель.}

\sectionline{3. НОК и НОД через простые множители}
\learn{Вы раскладываете числа на простые множители и записываете их в виде степеней.}
\learn{Для НОК берёте \emph{наибольшие} степени всех простых множителей,
для НОД --- \emph{наименьшие} степени.}

\textbf{Пример 4.} Найдём НОК и НОД чисел $6$, $16$, $50$:
\[
6=2\cdot3,\qquad 16=2^4,\qquad 50=2\cdot5^2.
\]
Для отсутствующего простого множителя удобно считать показатель степени
равным нулю ($3^0=5^0=1$).
\[
\text{НОК}(6,16,50)=2^4\cdot3\cdot5^2=\boxed{1200},
\]
\[
\text{НОД}(6,16,50)=2^1\cdot3^0\cdot5^0=\boxed{2}.
\]
\keyidea{НОК помогает найти общий знаменатель; НОД помогает сократить дробь.
В общий знаменатель входят все необходимые простые множители.}

\sectionline{4. Неизвестный компонент действия}
\learn{Вы узнаёте вид уравнения и применяете обратное действие.}
\learn{Вы решаете уравнение с дробями, затем проверяете подстановкой.}

\renewcommand{\arraystretch}{1.28}
\begin{center}
\begin{tabular}{@{}lll@{}}
\toprule
\textbf{Равенство} & \textbf{Найти $a$} & \textbf{Найти $b$}\\
\midrule
$a+b=c$ & $a=c-b$ & $b=c-a$\\
$a-b=c$ & $a=c+b$ & $b=a-c$\\
$a\cdot b=c$ & $a=c:b$ & $b=c:a$\\
$a:b=c$ & $a=b\cdot c$ & $b=a:c$\\
\bottomrule
\end{tabular}
\end{center}
\emph{Условие:} при делении делитель должен быть ненулевым.
Формулы с делением применяют, только когда указанный делитель не равен нулю.

\textbf{Пример 5.} Найдём неизвестное вычитаемое:
\begin{align*}
\mix{4}{1}{3}-x&=\mix{1}{5}{6},\\
x&=\mix{4}{1}{3}-\mix{1}{5}{6}
=\mix{4}{2}{6}-\mix{1}{5}{6}\\
&=3+\frac86-\mix{1}{5}{6}
=2+\frac36
=\boxed{\mix{2}{1}{2}}.
\end{align*}
\keyidea{В равенстве $a-x=b$ неизвестное вычитаемое равно $a-b$.
Не меняйте местами уменьшаемое и вычитаемое.}

\sectionline{5. Перевод слов на математический язык}
\learn{Вы различаете сумму, разность, произведение и частное.}
\learn{Вы ставите скобки так, чтобы сохранить смысл условия.}

\[
\text{сумма $a$ и $b$}=a+b,\qquad
\text{разность $a$ и $b$}=a-b.
\]
\textbf{Пример 6.} «Из суммы $\,\mix{1}{1}{4}\,$ и $\,\mix{2}{1}{6}\,$
вычесть $\,\mix{1}{1}{3}\,$» означает
\[
\left(\mix{1}{1}{4}+\mix{2}{1}{6}\right)-\mix{1}{1}{3}
=\mix{3}{5}{12}-\mix{1}{4}{12}
=\boxed{\mix{2}{1}{12}}.
\]

\textbf{Подсказка «ВУПС»:} \textbf{В}ычитаемое,
\textbf{У}множаемое, \textbf{П}одставляемое --- \textbf{С}кобки.
Когда составное выражение используется как одно число при вычитании
или умножении, скобки сохраняют нужный порядок действий.
Они нужны и вокруг суммы либо разности в делителе.

\sectionline{6. Текстовые задачи: три основные модели}
\learn{Вы выписываете данные и определяете величину, которую нужно найти.}
\learn{Вы проверяете единицы измерения и смысл результата.}

\textbf{Периметр.} Это сумма длин \emph{всех} сторон фигуры.
Для треугольника со сторонами $a$, $b$, $c$:
\[
P=a+b+c.
\]
Для четырёхугольника складывают четыре стороны и так далее.

\textbf{Урожайность.} Это масса собранного урожая с единицы площади:
\[
Y=\frac{m}{S},\qquad
\frac{\text{ц}}{\text{га}}=\text{центнеров с гектара}.
\]
Если собрали $1052$ ц с $25$ га и $1463$ ц с $30$ га, то
\[
Y_1=\frac{1052}{25}=\frac{6312}{150},\qquad
Y_2=\frac{1463}{30}=\frac{7315}{150}.
\]
Следовательно, урожайность второго поля выше на
\[
Y_2-Y_1=\frac{7315-6312}{150}
=\boxed{\mix{6}{103}{150}\ \text{ц/га}}.
\]

\textbf{Движение по реке.} Собственную скорость лодки находим по пути и времени:
\[
v_0=\frac{s}{t}.
\]
Скорость \emph{по течению} больше собственной на скорость течения,
а скорость \emph{против течения} --- меньше:
\[
v_{\text{по теч.}}=v_0+v_{\text{теч.}},\qquad
v_{\text{против теч.}}=v_0-v_{\text{теч.}}.
\]
Чтобы лодка продвигалась вверх по реке, должно быть $v_0>v_{\text{теч.}}$.
\[
s=v\cdot t,\qquad t=\frac{s}{v}\quad(v>0).
\]
\keyidea{Если спрашивают, \emph{насколько} одна величина больше другой,
найдите разность большего и меньшего. В ответе укажите единицы измерения.}

\clearpage
\thispagestyle{plain}
\sectionline{Алгоритм: от текста задачи к ответу}
\vspace{4mm}
\begin{center}
\begin{tikzpicture}[
  >=Stealth,
  every node/.style={align=center,font=\small},
  term/.style={draw=darkaccent,ellipse,minimum width=4cm,minimum height=1cm},
  action/.style={draw=accent,rounded corners=3pt,text width=9cm,
                 minimum height=1.18cm,inner sep=5pt},
  decision/.style={draw=darkaccent,diamond,aspect=2.5,
                   text width=4.2cm,inner sep=3pt},
  sideaction/.style={draw=accent,rounded corners=3pt,text width=3.2cm,
                     minimum height=1.7cm,inner sep=5pt},
  arrow/.style={->,line width=.9pt,draw=darkaccent},
  node distance=11mm
]
\node[term] (start) {Начало};
\node[action,below=of start] (data)
  {Выпишите известные величины, искомое и единицы измерения};
\node[action,below=of data] (model)
  {Выберите правило: сумма сторон, урожайность, движение или уравнение};
\node[action,below=of model] (calculate)
  {Составьте выражение; вычислите значения и сократите дроби};
\node[decision,below=16mm of calculate] (diff)
  {Нужно узнать,\\насколько больше?};
\node[sideaction,right=16mm of diff] (subtract)
  {Найдите разность: большее минус меньшее};
\node[action,below=21mm of diff] (check)
  {Проверьте порядок величин, единицы и смысл ответа};
\node[term,below=of check] (end) {Ответ};
\draw[arrow] (start)--(data);
\draw[arrow] (data)--(model);
\draw[arrow] (model)--(calculate);
\draw[arrow] (calculate)--(diff);
\draw[arrow] (diff.south)--node[left]{Нет}(check.north);
\draw[arrow] (diff.east)--node[above]{Да}(subtract.west);
\draw[arrow] (subtract.south) |- (check.east);
\draw[arrow] (check)--(end);
\end{tikzpicture}
\end{center}

\clearpage
\sectionline{7. Самостоятельная тренировка (Путь А)}
\textbf{10 заданий.} Решайте последовательно. Для вычислений запишите основные
преобразования; для текстовых задач --- выражение, вычисления и ответ с единицами.
В задании 4 покажите разложение на простые множители.

\begin{enumerate}[itemsep=1.08em]
\item Вычислите удобным способом:
\[
\mix{1}{2}{9}+\mix{2}{5}{12}+\mix{3}{7}{9}.
\]
\item Раскройте скобки и вычислите:
\[
\mix{6}{3}{8}-\left(\mix{2}{1}{8}+\mix{1}{5}{8}\right).
\]
\item Представьте десятичную дробь в виде смешанного числа,
сократите дробную часть и вычислите:
\[
\mix{3}{3}{10}-1{,}84+\mix{2}{1}{4}.
\]
\item Найдите $\text{НОД}(12,18,45)$ и
$\text{НОК}(12,18,45)$ методом простых множителей.
\item Решите уравнение:
\[
\mix{5}{1}{6}-x=\mix{2}{3}{4}.
\]
\item Составьте выражение и вычислите его значение:
«Из суммы чисел $\,\mix{2}{1}{3}\,$ и $\,\mix{1}{1}{4}\,$
вычесть число $\,\mix{1}{5}{6}\,$».
\item Найдите периметр треугольника со сторонами
$\,\mix{4}{1}{6}\,$ см, $\,\mix{5}{3}{8}\,$ см и $\,\mix{3}{1}{4}\,$ см.
\item С первого поля площадью $15$ га собрали $635$ ц пшеницы,
а со второго поля площадью $18$ га --- $754$ ц.
На каком поле урожайность выше и на сколько центнеров с гектара?
\item Моторная лодка в стоячей воде проходит $1500$ м за $8$ мин.
Скорость течения реки $\,\mix{7}{3}{4}\,$ м/мин.
Найдите скорости лодки по течению и против течения.
\item Из суммы неизвестного числа и $\,\mix{2}{1}{6}\,$
вычли $\,\mix{1}{3}{4}\,$ и получили $\,\mix{4}{1}{8}\,$.
Составьте уравнение и найдите неизвестное число.
\end{enumerate}

\clearpage
\sectionline{8. Ответы для самопроверки}
\renewcommand{\arraystretch}{1.4}
\begin{tabularx}{\linewidth}{@{}p{10mm}X@{}}
\toprule
\textbf{№} & \textbf{Ответ и ключевой шаг}\\
\midrule
1 & $\mix{7}{5}{12}$.
Сначала $\frac29+\frac79=1$.\\
2 & $\mix{2}{5}{8}$.
В скобках $\mix{3}{3}{4}$.\\
3 & $\mix{3}{71}{100}$.
$1{,}84=\mix{1}{21}{25}$, общий знаменатель $100$.\\
4 & $\text{НОД}=3$, $\text{НОК}=180$;
$12=2^2\cdot3$, $18=2\cdot3^2$, $45=3^2\cdot5$.\\
5 & $x=\mix{2}{5}{12}$.
$x=\mix{5}{1}{6}-\mix{2}{3}{4}$.\\
6 & $\mix{1}{3}{4}$.
$(\mix{2}{1}{3}+\mix{1}{1}{4})-\mix{1}{5}{6}$.\\
7 & $\mix{12}{19}{24}$ см.
Общий знаменатель дробных частей $24$.\\
8 & На первом поле; на $\frac49$ ц/га.
$Y_1=\mix{42}{1}{3}$ ц/га, $Y_2=\mix{41}{8}{9}$ ц/га.\\
9 & По течению $\mix{195}{1}{4}$ м/мин,
против течения $\mix{179}{3}{4}$ м/мин;
$v_0=\frac{1500}{8}=\mix{187}{1}{2}$ м/мин.\\
10 & $x=\mix{3}{17}{24}$.
$(x+\mix{2}{1}{6})-\mix{1}{3}{4}=\mix{4}{1}{8}$.\\
\bottomrule
\end{tabularx}
\vspace{8mm}
\keyidea{Сверяйте способ решения и ответ. Если результат отличается,
сначала проверьте знаки перед скобками, общий знаменатель и «заём».}
\end{document}
